\section{数据特征与证据分层}

\subsection{时间连续性、缺失与疑似量程封顶}

原始数据从 2024 年 5 月 1 日 00:00 连续至 5 月 7 日 23:59，共 $7\times24\times60=10,080$ 行，时间步长均为 1 分钟。除出口浓度缺失 50 行外，其余建模变量完整。图\ref{fig:dataaudit} 给出数据时序与出口边界分布。

\begin{figure}[H]
\centering
\includegraphics[width=0.97\textwidth]{fig01_data_audit.png}
\caption{七天运行数据的时序与出口浓度边界堆积}
\label{fig:dataaudit}
\end{figure}

在 10,030 个有效出口读数中，5,491 个恰好等于 $50\mgNm$，比例为 54.75\%；若以含缺失值的 10,080 个时刻为分母则为 54.47\%。同时，数据中不存在高于 50 的读数。连续分布与边界点集中并存的特征表明，仪表或数据接口可能存在上限截顶，但题面未直接给出量程说明。因此，本文将上限删失作为建模假设，定义观测变量
\begin{equation}
Y_t=\min(Y_t^*,c),\qquad c=50,
\label{eq:censorobs}
\end{equation}
其中 $Y_t^*$ 为潜在真实出口浓度。该式是依据数据分布特征提出的工作假设，并非题目给定条件；缺失值仅在出口浓度建模时剔除，不进行插值，以避免引入额外的排放轨迹假设。

\begin{table}[H]
\centering
\caption{数据质量检验结果}
\label{tab:audit}
\begin{tabularx}{0.92\textwidth}{L C{3.0cm} L}
\toprule
检验项目 & 结果 & 建模处理\\
\midrule
时间长度 & 10,080 min & 保持原始顺序，验证按天分组\\
时间间隔 & 全部为 1 min & 保持原有时间网格\\
出口缺失 & 50 行 & 不插补，仅在排放似然中剔除\\
疑似出口上限 & 5,491 行，占有效读数 54.75\% & 基准情景按 $Y^*\ge50$ 的删失事件处理\\
原始出口普通 OLS & $R^2=0.000815$ & 不用于因果解释或限值外推\\
封顶事件线性模型 & $R^2=0.001162$ & 条件依赖不可识别\\
功率代理全样本 & $R^2=0.9979$ & 可用于历史边界附近的优化\\
\bottomrule
\end{tabularx}
\end{table}

\subsection{删失正态运行点估计}

为估计量程附近的潜在中心和噪声，令 $Y_t^*\sim\mathcal N(\mu_c,\sigma_c^2)$。记未删失集合为 $\mathcal U$、封顶集合为 $\mathcal C$，删失对数似然为
\begin{equation}
\ell(\mu_c,\sigma_c)=
\sum_{t\in\mathcal U}\left[-\log\sigma_c+\log\phi\!\left(\frac{Y_t-\mu_c}{\sigma_c}\right)\right]
+\sum_{t\in\mathcal C}\log\!\left[1-\Phi\!\left(\frac{c-\mu_c}{\sigma_c}\right)\right].
\label{eq:censorll}
\end{equation}
极大似然估计为 $\hat\mu_c=50.036\mgNm$、$\hat\sigma_c=0.313\mgNm$。该结果仅用于刻画观测期内的闭环运行水平，不足以辨识电压、温度或振打周期的净效应。图\ref{fig:censordiag} 中的相关系数接近零，进一步支持这一判断。

\begin{figure}[H]
\centering
\includegraphics[width=0.91\textwidth]{fig02_censor_diagnostics.png}
\caption{出口浓度相关性与删失假设下的运行点估计}
\label{fig:censordiag}
\end{figure}

\subsection{闭环偏差与证据等级}

生产控制器会在高入口负荷时主动提高电压或改变振打周期，导致“高电压与高出口浓度同时出现”的表观相关。此类闭环偏差意味着相关关系不能直接解释为控制量导致排放上升。更关键的是，54.75\% 的有效读数封顶使输出变化严重压缩。故本文采用表\ref{tab:evidence} 的证据分层：数据用于确定运行点、能耗和工况分布；机理参数用于限值外推，并对其作显式敏感性分析。

\begin{table}[H]
\centering
\caption{排放模型参数的证据等级}
\label{tab:evidence}
\begin{tabularx}{0.96\textwidth}{L C{2.6cm} C{3.8cm} C{2.2cm}}
\toprule
对象 & 证据来源 & 诊断量 & 结论\\
\midrule
边界位置与局部噪声 & 附件数据+删失假设 & $n=10,030$，边界占比 54.746\% & 条件可识别\\
工况/操作量对出口的净效应 & 附件数据 & 普通 OLS $R^2=0.000815$ & 不可识别\\
删失事件的条件依赖 & 附件数据 & 线性模型 $R^2=0.001162$ & 不可识别\\
基准等效去除指数 & 删失锚定+质量守恒 & $D_0=6.6846$ & 弱识别\\
分场权重 & 附件中位电压+D--A平方代理 & $w_i=U_{i,0}^2/\sum_jU_{j,0}^2$ & 代理识别\\
温度参考中心 & 附件数据 & 均值 $126.008^\circ\mathrm C$ & 可识别\\
温度曲率、积灰与再飞扬 & 公开机理研究 & 基准情景传播，并作 $\pm30\%$ 敏感性分析 & 先验主导\\
\bottomrule
\end{tabularx}
\end{table}

\begin{warnbox}
\textbf{外推范围：} 附件的有效出口数据集中在 $50\mgNm$ 附近，而题目要求 $10\mgNm$ 与 $5\mgNm$，二者均位于直接观测范围之外。因此，后续浓度结果属于删失运行点与物理先验共同约束下的情景外推，尚需低量程现场数据进一步验证。
\end{warnbox}
